% GENERATED FILE. Edit canonical lessons, then run npm run generate:books.
% canonical-source-hash: fnv1a64:a214a859dc46a357
% canonical-lessons: TE-C453-oddu, TE-C453-vagu, TE-C453-teppa, TE-C453-telipovu, TE-C453-kottukupovu, TE-R453-review-446-453

\chapter{On the Water}
\label{ch:te-on-the-water}
\hlchaptermodality{453}

\begin{chapteropening}
\textbf{By the end of this chapter:} \emph{I can say a riverbank, a rivulet, a raft, to float away and to be swept away.}

Complete the last lesson of chapter 453: I can say a riverbank, a rivulet, a raft, to float away and to be swept away.
\end{chapteropening}

\section[\texorpdfstring{\te{ఒడ్డు} (oḍḍu)}{ఒడ్డు (oḍḍu)}]{\te{ఒడ్డు} (oḍḍu) — a riverbank; a shore}

\label{lesson:TE-C453-oddu}



\begin{quote}
Before the new one: say the Telugu for an anchor, then the Telugu for a lighthouse.
\end{quote}

\subsection*{You'll want to know: \te{ఒడ్డు}}
\textbf{\te{ఒడ్డు}} — \emph{oḍḍu} — \textquotedblleft{}a riverbank; a shore\textquotedblright{}.

\begin{quote}
\te{పిల్లలు} \te{నది} \te{ఒడ్డున} \te{ఆడుకుంటున్నారు}.
\end{quote}

\begin{quote}
\emph{pillalu nadi oḍḍuna āḍukuṇṭunnāru.}
\end{quote}

\begin{quote}
The children are playing on the riverbank.
\end{quote}

\te{ఒడ్డున}, on the bank, is the usual form. The bank of a pond, a river or a sea.

\begin{cousinweb}
\textbf{\te{ఒడ్డు}} (\emph{oḍḍu}) is native Dravidian: DEDR 965 lists Telugu \emph{oḍḍu}, \textquotedblleft{}a bank\textquotedblright{}, with Kannada \emph{oḍḍu}, \textquotedblleft{}a bank, a causeway\textquotedblright{}.
\end{cousinweb}

\begin{grammarlens}[title={What you've built}]
One more word for getting around.
\end{grammarlens}

\subsection*{Your turn}
\begin{itemize}
\raggedright
  \item \emph{Say it:} \emph{oḍḍu}
  \item \emph{Say it:} \emph{oḍḍu}, once more
  \item \emph{Say it:} \emph{pillalu nadi oḍḍuna āḍukuṇṭunnāru.} — The children are playing on the riverbank.
  \item \emph{Recall:} say the Telugu for an anchor, then the Telugu for a lighthouse, then say \emph{oḍḍu} again
\end{itemize}

\begin{tcolorbox}[breakable,colback=teal!4,colframe=teal!35!black,title={Before you move on}]
What does \te{ఒడ్డు} mean? (\textquotedblleft{}A riverbank; a shore\textquotedblright{}.) Say it once more.
\end{tcolorbox}
\section[\texorpdfstring{\te{వాగు} (vāgu)}{వాగు (vāgu)}]{\te{వాగు} (vāgu) — a rivulet; a stream}

\label{lesson:TE-C453-vagu}



\begin{quote}
Before the new one: say the Telugu for a lighthouse, then the Telugu for a riverbank.
\end{quote}

\subsection*{You'll want to know: \te{వాగు}}
\textbf{\te{వాగు}} — \emph{vāgu} — \textquotedblleft{}a rivulet; a stream\textquotedblright{}.

\begin{quote}
\te{వర్షాకాలంలో} \te{ఆ} \te{వాగు} \te{నిండుగా} \te{పారుతుంది}.
\end{quote}

\begin{quote}
\emph{varṣākālaṁlō ā vāgu niṇḍugā pārutundi.}
\end{quote}

\begin{quote}
In the rainy season that stream runs full.
\end{quote}

A small stream that fills in the rains and dries in summer.

\begin{cousinweb}
\textbf{\te{వాగు}} (\emph{vāgu}) is native Dravidian: DEDR 5334 lists Telugu \emph{vāgu}, \textquotedblleft{}a rivulet, a streamlet\textquotedblright{}, in the family of Tamil \emph{vākku}, \textquotedblleft{}to pour\textquotedblright{}.
\end{cousinweb}

\begin{grammarlens}[title={What you've built}]
One more word for getting around.
\end{grammarlens}

\subsection*{Your turn}
\begin{itemize}
\raggedright
  \item \emph{Say it:} \emph{vāgu}
  \item \emph{Say it:} \emph{vāgu}, once more
  \item \emph{Say it:} \emph{varṣākālaṁlō ā vāgu niṇḍugā pārutundi.} — In the rainy season that stream runs full.
  \item \emph{Recall:} say the Telugu for a lighthouse, then the Telugu for a riverbank, then say \emph{vāgu} again
\end{itemize}

\begin{tcolorbox}[breakable,colback=teal!4,colframe=teal!35!black,title={Before you move on}]
What does \te{వాగు} mean? (\textquotedblleft{}A rivulet; a stream\textquotedblright{}.) Say it once more.
\end{tcolorbox}
\section[\texorpdfstring{\te{తెప్ప} (teppa)}{తెప్ప (teppa)}]{\te{తెప్ప} (teppa) — a raft}

\label{lesson:TE-C453-teppa}



\begin{quote}
Before the new one: say the Telugu for a riverbank, then the Telugu for a rivulet.
\end{quote}

\subsection*{You'll want to know: \te{తెప్ప}}
\textbf{\te{తెప్ప}} — \emph{teppa} — \textquotedblleft{}a raft\textquotedblright{}.

\begin{quote}
\te{వాళ్ళు} \te{తెప్ప} \te{మీద} \te{వాగు} \te{దాటారు}.
\end{quote}

\begin{quote}
\emph{vāḷḷu teppa mīda vāgu dāṭāru.}
\end{quote}

\begin{quote}
They crossed the stream on a raft.
\end{quote}

A raft of logs or bamboo tied together.

\begin{cousinweb}
\textbf{\te{తెప్ప}} (\emph{teppa}) is native Dravidian: DEDR 3414 lists Telugu \emph{teppa}, \textquotedblleft{}a raft, a catamaran\textquotedblright{}, with Tamil \emph{teppam}.
\end{cousinweb}

\begin{grammarlens}[title={What you've built}]
One more word for getting around.
\end{grammarlens}

\subsection*{Your turn}
\begin{itemize}
\raggedright
  \item \emph{Say it:} \emph{teppa}
  \item \emph{Say it:} \emph{teppa}, once more
  \item \emph{Say it:} \emph{vāḷḷu teppa mīda vāgu dāṭāru.} — They crossed the stream on a raft.
  \item \emph{Recall:} say the Telugu for a riverbank, then the Telugu for a rivulet, then say \emph{teppa} again
\end{itemize}

\begin{tcolorbox}[breakable,colback=teal!4,colframe=teal!35!black,title={Before you move on}]
What does \te{తెప్ప} mean? (\textquotedblleft{}A raft\textquotedblright{}.) Say it once more.
\end{tcolorbox}
\section[\texorpdfstring{\te{తేలిపోవు} (tēlipōvu)}{తేలిపోవు (tēlipōvu)}]{\te{తేలిపోవు} (tēlipōvu) — to float away}

\label{lesson:TE-C453-telipovu}



\begin{quote}
Before the new one: say the Telugu for a rivulet, then the Telugu for a raft.
\end{quote}

\subsection*{You'll want to know: \te{తేలిపోవు}}
\textbf{\te{తేలిపోవు}} — \emph{tēlipōvu} — \textquotedblleft{}to float away\textquotedblright{}.

\begin{quote}
\te{గాలికి} \te{నా} \te{టోపీ} \te{నీటిలో} \te{తేలిపోయింది}.
\end{quote}

\begin{quote}
\emph{gāliki nā ṭōpī nīṭilō tēlipōyindi.}
\end{quote}

\begin{quote}
In the wind my cap floated away on the water.
\end{quote}

Past (it …): \te{తేలిపోయింది}, \emph{tēlipōyindi}.

\te{తేలు}, to float, + \te{పోవు}, to go: carried off on the surface. Its subject is a thing, so the past shown is the third person: \te{తేలిపోయింది}, it floated away.

\begin{cousinweb}
\textbf{\te{తేలిపోవు}} (\emph{tēlipōvu}) is native Dravidian: \te{తేలు}, \textquotedblleft{}to float\textquotedblright{} (DEDR 3464), with \te{పోవు}, \textquotedblleft{}to go\textquotedblright{} (DEDR 4572).
\end{cousinweb}

\begin{grammarlens}[title={What you've built}]
One more word for getting around.
\end{grammarlens}

\subsection*{Your turn}
\begin{itemize}
\raggedright
  \item \emph{Say it:} \emph{tēlipōvu}
  \item \emph{Say it:} \emph{tēlipōvu}, once more
  \item \emph{Say it:} \emph{gāliki nā ṭōpī nīṭilō tēlipōyindi.} — In the wind my cap floated away on the water.
  \item \emph{Recall:} say the Telugu for a rivulet, then the Telugu for a raft, then say \emph{tēlipōvu} again
\end{itemize}

\begin{tcolorbox}[breakable,colback=teal!4,colframe=teal!35!black,title={Before you move on}]
What does \te{తేలిపోవు} mean? (\textquotedblleft{}To float away\textquotedblright{}.) Say it once more.
\end{tcolorbox}
\section[\texorpdfstring{\te{కొట్టుకుపోవు} (koṭṭukupōvu)}{కొట్టుకుపోవు (koṭṭukupōvu)}]{\te{కొట్టుకుపోవు} (koṭṭukupōvu) — to be swept away}

\label{lesson:TE-C453-kottukupovu}



\begin{quote}
Before the new one: say the Telugu for a raft, then the Telugu for to float away.
\end{quote}

\subsection*{You'll want to know: \te{కొట్టుకుపోవు}}
\textbf{\te{కొట్టుకుపోవు}} — \emph{koṭṭukupōvu} — \textquotedblleft{}to be swept away\textquotedblright{}.

\begin{quote}
\te{వరదలో} \te{చాలా} \te{పశువులు} \te{కొట్టుకుపోయాయి}.
\end{quote}

\begin{quote}
\emph{varadalō cālā paśuvulu koṭṭukupōyāyi.}
\end{quote}

\begin{quote}
Many cattle were swept away in the flood.
\end{quote}

Past (it …): \te{కొట్టుకుపోయింది}, \emph{koṭṭukupōyindi}.

Said of anything a current carries off: a bridge, a road, a crop. Its subject is a thing, so the past shown is the third person: \te{కొట్టుకుపోయింది}, it was swept away.

\begin{cousinweb}
\textbf{\te{కొట్టుకుపోవు}} (\emph{koṭṭukupōvu}) is native Dravidian: DEDR 2065 lists Telugu \emph{koṭṭukonipōvu}, \textquotedblleft{}to be carried along by a stream or a current\textquotedblright{}, from \emph{koṭṭu}, \textquotedblleft{}to pour forth\textquotedblright{}.
\end{cousinweb}

\begin{grammarlens}[title={What you've built}]
One more word for getting around.
\end{grammarlens}

\subsection*{Your turn}
\begin{itemize}
\raggedright
  \item \emph{Say it:} \emph{koṭṭukupōvu}
  \item \emph{Say it:} \emph{koṭṭukupōvu}, once more
  \item \emph{Say it:} \emph{varadalō cālā paśuvulu koṭṭukupōyāyi.} — Many cattle were swept away in the flood.
  \item \emph{Recall:} say the Telugu for a raft, then the Telugu for to float away, then say \emph{koṭṭukupōvu} again
\end{itemize}

\begin{tcolorbox}[breakable,colback=teal!4,colframe=teal!35!black,title={Before you move on}]
What does \te{కొట్టుకుపోవు} mean? (\textquotedblleft{}To be swept away\textquotedblright{}.) Say it once more.
\end{tcolorbox}
\section[(dialogue)]{Words from chapters 446–453 — a review}

\label{lesson:TE-R453-review-446-453}



\begin{quote}
Say the Telugu for bumpy and a jolt.
\end{quote}

\subsection*{Your turn}
Say these aloud:

\begin{itemize}
\raggedright
  \item bumpy, a jolt, to avoid, to overturn, to get dented
  \item a traffic signal, a licence, a pedestrian, a footpath, to break a rule
  \item a compass direction, a signpost, to lose one's way, to take a turn, a shortcut
  \item an obstacle, a diversion, to stop dead, to seize up, to get blocked
  \item parking, to pull up, to put aboard, to unload, to drop someone off
  \item a rickshaw, a bullock cart, a handcart, going on foot, a wayfarer
  \item a harbour, a ship, a sailor, an anchor, a lighthouse
  \item a riverbank, a rivulet, a raft, to float away, to be swept away
\end{itemize}

\begin{tcolorbox}[breakable,colback=teal!4,colframe=teal!35!black,title={Before you move on}]
Bumpy? (\textbf{\te{ఎగుడుదిగుడు}}, \emph{eguḍudiguḍu}.) A jolt? (\textbf{\te{కుదుపు}}, \emph{kudupu}.) To avoid? (\textbf{\te{తప్పించు}}, \emph{tappiñcu}.) To overturn? (\textbf{\te{బోల్తాపడు}}, \emph{bōltāpaḍu}.) To get dented? (\textbf{\te{సొట్టపడు}}, \emph{soṭṭapaḍu}.) A traffic signal? (\textbf{\te{సిగ్నలు}}, \emph{signalu}.) A licence? (\textbf{\te{లైసెన్సు}}, \emph{laisensu}.) A pedestrian? (\textbf{\te{పాదచారి}}, \emph{pādacāri}.) A footpath? (\textbf{\te{కాలిబాట}}, \emph{kālibāṭa}.) To break a rule? (\textbf{\te{అతిక్రమించు}}, \emph{atikramiñcu}.) A compass direction? (\textbf{\te{దిశ}}, \emph{diśa}.) A signpost? (\textbf{\te{సూచిక}}, \emph{sūcika}.) To lose one's way? (\textbf{\te{దారి} \te{తప్పు}}, \emph{dāri tappu}.) To take a turn? (\textbf{\te{మలుపు} \te{తిరుగు}}, \emph{malupu tirugu}.) A shortcut? (\textbf{\te{అడ్డదారి}}, \emph{aḍḍadāri}.) An obstacle? (\textbf{\te{అడ్డంకి}}, \emph{aḍḍaṅki}.) A diversion? (\textbf{\te{మళ్ళింపు}}, \emph{maḷḷimpu}.) To stop dead? (\textbf{\te{నిలిచిపోవు}}, \emph{nilicipōvu}.) To seize up? (\textbf{\te{స్తంభించు}}, \emph{stambhiñcu}.) To get blocked? (\textbf{\te{మూసుకుపోవు}}, \emph{mūsukupōvu}.) Parking? (\textbf{\te{పార్కింగు}}, \emph{pārkiṅgu}.) To pull up? (\textbf{\te{నిలుపు}}, \emph{nilupu}.) To put aboard? (\textbf{\te{ఎక్కించు}}, \emph{ekkiñcu}.) To unload? (\textbf{\te{దించు}}, \emph{diñcu}.) To drop someone off? (\textbf{\te{దిగబెట్టు}}, \emph{digabeṭṭu}.) A rickshaw? (\textbf{\te{రిక్షా}}, \emph{rikṣā}.) A bullock cart? (\textbf{\te{ఎడ్లబండి}}, \emph{eḍlabaṇḍi}.) A handcart? (\textbf{\te{తోపుడుబండి}}, \emph{tōpuḍubaṇḍi}.) Going on foot? (\textbf{\te{కాలినడక}}, \emph{kālinaḍaka}.) A wayfarer? (\textbf{\te{బాటసారి}}, \emph{bāṭasāri}.) A harbour? (\textbf{\te{రేవు}}, \emph{rēvu}.) A ship? (\textbf{\te{నౌక}}, \emph{nauka}.) A sailor? (\textbf{\te{నావికుడు}}, \emph{nāvikuḍu}.) An anchor? (\textbf{\te{లంగరు}}, \emph{laṅgaru}.) A lighthouse? (\textbf{\te{దీపస్తంభం}}, \emph{dīpastambhaṁ}.) A riverbank? (\textbf{\te{ఒడ్డు}}, \emph{oḍḍu}.) A rivulet? (\textbf{\te{వాగు}}, \emph{vāgu}.) A raft? (\textbf{\te{తెప్ప}}, \emph{teppa}.) To float away? (\textbf{\te{తేలిపోవు}}, \emph{tēlipōvu}.) To be swept away? (\textbf{\te{కొట్టుకుపోవు}}, \emph{koṭṭukupōvu}.)
\end{tcolorbox}
